\chapter{Messdaten und Reproduzierbarkeit}
\label{chap:anhang_messdaten}

Dieser Anhang enthält die Messwerte hinter Kapitel~\ref{chap:ergebnisse} und ordnet jede Abbildung und Tabelle dem Skript zu, das sie erzeugt oder prüft. Die Tabellen~\ref{tab:anh_prf1} bis~\ref{tab:anh_zuordnung} schreibt \texttt{anhang\_\allowbreak tabellen.py} direkt aus den CSV-Dateien, keine Zahl ist von Hand übertragen.

\section{Messwerte aller Verfahren}
\label{sec:anh_messwerte}

Tabelle~\ref{tab:anh_prf1} fasst Precision, Recall und F1-Score aller vier Verfahren zusammen. Die F1-Spalte stimmt mit Tabelle~\ref{tab:f1_vergleich} überein. Precision und Recall der Baselines erscheinen im Hauptteil vor allem als Arbeitspunkte in Abbildung~\ref{fig:pr_ebene}. Tabelle~\ref{tab:anh_seeds} listet die Einzelwerte, aus denen Mittelwert und Standardabweichung gebildet sind. Die Standardabweichung folgt der Populationsformel mit $n=3$, so wie sie \texttt{train\_\allowbreak eval.py} und \texttt{eval\_\allowbreak baselines.py} berechnen.

\begin{table}[htbp]
	\centering
	\small
	\renewcommand{\arraystretch}{1.15}
	\caption[Precision, Recall und F1-Score aller Verfahren je SNR-Stufe]{Precision, Recall und F1-Score aller Verfahren je SNR-Stufe, Mittelwert $\pm$ Standardabweichung über die Seeds 0, 1 und 2 mit je $20\,000$ Testflüssen.}
	\label{tab:anh_prf1}
	% erzeugt von anhang_tabellen.py, nicht von Hand ändern
\begin{tabular}{llccc}
\toprule
\textbf{Verfahren} & \textbf{SNR in dB} & \textbf{Precision} & \textbf{Recall} & \textbf{F1-Score} \\
\midrule
DDQN & 30 & $0{,}8853 \pm 0{,}0021$ & $0{,}9018 \pm 0{,}0022$ & $0{,}8935 \pm 0{,}0009$ \\
 & 20 & $0{,}8813 \pm 0{,}0023$ & $0{,}8977 \pm 0{,}0015$ & $0{,}8894 \pm 0{,}0016$ \\
 & 10 & $0{,}8120 \pm 0{,}0048$ & $0{,}8485 \pm 0{,}0029$ & $0{,}8298 \pm 0{,}0037$ \\
 & 5 & $0{,}6963 \pm 0{,}0042$ & $0{,}7583 \pm 0{,}0031$ & $0{,}7260 \pm 0{,}0037$ \\
\midrule
Random Forest & 30 & $0{,}9260 \pm 0{,}0036$ & $0{,}8274 \pm 0{,}0076$ & $0{,}8739 \pm 0{,}0032$ \\
 & 20 & $0{,}9177 \pm 0{,}0046$ & $0{,}8202 \pm 0{,}0064$ & $0{,}8662 \pm 0{,}0026$ \\
 & 10 & $0{,}8551 \pm 0{,}0069$ & $0{,}7582 \pm 0{,}0102$ & $0{,}8037 \pm 0{,}0031$ \\
 & 5 & $0{,}7327 \pm 0{,}0067$ & $0{,}6649 \pm 0{,}0067$ & $0{,}6972 \pm 0{,}0060$ \\
\midrule
SVM (RBF) & 30 & $0{,}9236 \pm 0{,}0024$ & $0{,}8508 \pm 0{,}0045$ & $0{,}8857 \pm 0{,}0020$ \\
 & 20 & $0{,}9183 \pm 0{,}0030$ & $0{,}8439 \pm 0{,}0025$ & $0{,}8795 \pm 0{,}0013$ \\
 & 10 & $0{,}8572 \pm 0{,}0068$ & $0{,}7880 \pm 0{,}0055$ & $0{,}8211 \pm 0{,}0016$ \\
 & 5 & $0{,}7353 \pm 0{,}0026$ & $0{,}6858 \pm 0{,}0066$ & $0{,}7097 \pm 0{,}0047$ \\
\midrule
MLP & 30 & $0{,}9099 \pm 0{,}0082$ & $0{,}8730 \pm 0{,}0079$ & $0{,}8910 \pm 0{,}0010$ \\
 & 20 & $0{,}9017 \pm 0{,}0115$ & $0{,}8675 \pm 0{,}0133$ & $0{,}8841 \pm 0{,}0022$ \\
 & 10 & $0{,}8356 \pm 0{,}0148$ & $0{,}8180 \pm 0{,}0142$ & $0{,}8265 \pm 0{,}0064$ \\
 & 5 & $0{,}7093 \pm 0{,}0168$ & $0{,}7315 \pm 0{,}0179$ & $0{,}7198 \pm 0{,}0003$ \\
\bottomrule
\end{tabular}

\end{table}

\begin{table}[tb]
	\centering
	\small
	\setlength{\tabcolsep}{3.5pt}
	\renewcommand{\arraystretch}{1.15}
	\caption[Einzelwerte je Seed]{Precision (P), Recall (R) und F1-Score je Seed und SNR-Stufe (SNR in dB).}
	\label{tab:anh_seeds}
	% erzeugt von anhang_tabellen.py, nicht von Hand ändern
\begin{tabular}{@{}lc*{9}{c}@{}}
\toprule
 & & \multicolumn{3}{c}{\textbf{Seed 0}} & \multicolumn{3}{c}{\textbf{Seed 1}} & \multicolumn{3}{c}{\textbf{Seed 2}} \\
\cmidrule(lr){3-5}\cmidrule(lr){6-8}\cmidrule(lr){9-11}
\textbf{Verfahren} & \textbf{SNR} & P & R & F1 & P & R & F1 & P & R & F1 \\
\midrule
DDQN & 30 & $0{,}8826$ & $0{,}9027$ & $0{,}8925$ & $0{,}8877$ & $0{,}8987$ & $0{,}8932$ & $0{,}8856$ & $0{,}9040$ & $0{,}8947$ \\
 & 20 & $0{,}8815$ & $0{,}8958$ & $0{,}8886$ & $0{,}8840$ & $0{,}8994$ & $0{,}8916$ & $0{,}8784$ & $0{,}8979$ & $0{,}8880$ \\
 & 10 & $0{,}8118$ & $0{,}8507$ & $0{,}8308$ & $0{,}8179$ & $0{,}8503$ & $0{,}8338$ & $0{,}8062$ & $0{,}8444$ & $0{,}8248$ \\
 & 5 & $0{,}7004$ & $0{,}7598$ & $0{,}7289$ & $0{,}6905$ & $0{,}7540$ & $0{,}7208$ & $0{,}6981$ & $0{,}7611$ & $0{,}7283$ \\
\midrule
RF & 30 & $0{,}9237$ & $0{,}8374$ & $0{,}8785$ & $0{,}9310$ & $0{,}8192$ & $0{,}8715$ & $0{,}9233$ & $0{,}8256$ & $0{,}8717$ \\
 & 20 & $0{,}9163$ & $0{,}8280$ & $0{,}8699$ & $0{,}9239$ & $0{,}8123$ & $0{,}8645$ & $0{,}9130$ & $0{,}8202$ & $0{,}8641$ \\
 & 10 & $0{,}8481$ & $0{,}7717$ & $0{,}8081$ & $0{,}8646$ & $0{,}7469$ & $0{,}8015$ & $0{,}8525$ & $0{,}7561$ & $0{,}8014$ \\
 & 5 & $0{,}7280$ & $0{,}6688$ & $0{,}6971$ & $0{,}7279$ & $0{,}6555$ & $0{,}6898$ & $0{,}7422$ & $0{,}6705$ & $0{,}7046$ \\
\midrule
SVM & 30 & $0{,}9232$ & $0{,}8563$ & $0{,}8885$ & $0{,}9266$ & $0{,}8452$ & $0{,}8840$ & $0{,}9208$ & $0{,}8509$ & $0{,}8845$ \\
 & 20 & $0{,}9174$ & $0{,}8472$ & $0{,}8809$ & $0{,}9223$ & $0{,}8412$ & $0{,}8799$ & $0{,}9152$ & $0{,}8434$ & $0{,}8778$ \\
 & 10 & $0{,}8538$ & $0{,}7942$ & $0{,}8229$ & $0{,}8667$ & $0{,}7809$ & $0{,}8215$ & $0{,}8512$ & $0{,}7891$ & $0{,}8190$ \\
 & 5 & $0{,}7371$ & $0{,}6883$ & $0{,}7119$ & $0{,}7317$ & $0{,}6767$ & $0{,}7031$ & $0{,}7371$ & $0{,}6923$ & $0{,}7140$ \\
\midrule
MLP & 30 & $0{,}9207$ & $0{,}8620$ & $0{,}8904$ & $0{,}9083$ & $0{,}8772$ & $0{,}8925$ & $0{,}9008$ & $0{,}8799$ & $0{,}8902$ \\
 & 20 & $0{,}9167$ & $0{,}8489$ & $0{,}8815$ & $0{,}8996$ & $0{,}8745$ & $0{,}8869$ & $0{,}8889$ & $0{,}8790$ & $0{,}8839$ \\
 & 10 & $0{,}8515$ & $0{,}7984$ & $0{,}8241$ & $0{,}8395$ & $0{,}8311$ & $0{,}8353$ & $0{,}8159$ & $0{,}8247$ & $0{,}8203$ \\
 & 5 & $0{,}7330$ & $0{,}7063$ & $0{,}7194$ & $0{,}6988$ & $0{,}7420$ & $0{,}7197$ & $0{,}6960$ & $0{,}7462$ & $0{,}7202$ \\
\bottomrule
\end{tabular}

\end{table}

\section{Implizite Schwellen}
\label{sec:anh_schwellen}

Die implizite Schwelle $t$ eines Modells ist die Posterior-Schwelle, bei der der Bayes-Klassifikator mit bekanntem SNR denselben Recall erreicht wie das Modell (Abschnitt~\ref{sec:reward_recall}). Tabelle~\ref{tab:anh_schwellen} gibt sie dreistellig an, berechnet aus den mittleren Recall-Werten der Tabelle~\ref{tab:anh_prf1}. Kapitel~\ref{chap:ergebnisse} rundet auf zwei Stellen. Abstände zwischen zwei Schwellen sind aus den ungerundeten Werten gebildet und lassen sich mit dieser Tabelle nachrechnen.

\begin{table}[htbp]
	\centering
	\small
	\renewcommand{\arraystretch}{1.15}
	\caption[Implizite Schwellen aller Verfahren]{Implizite Posterior-Schwelle $t$ je Verfahren und SNR-Stufe ($\pi=0{,}3$) und Soll-Schwelle des Lernziels. Mit Huber-Verlust liegt die Nullstelle des DDQN bei $3/7$ statt $1/3$ (Abschnitt~\ref{sec:dqn_architektur}).}
	\label{tab:anh_schwellen}
	% erzeugt von anhang_tabellen.py, nicht von Hand ändern
\begin{tabular}{lccccc}
\toprule
 & \multicolumn{4}{c}{\textbf{Implizite Schwelle $t$}} & \\
\cmidrule(lr){2-5}
\textbf{Verfahren} & 30\,dB & 20\,dB & 10\,dB & 5\,dB & \textbf{Soll-Schwelle} \\
\midrule
DDQN & $0{,}428$ & $0{,}415$ & $0{,}382$ & $0{,}375$ & $1/3$ \\
Random Forest & $0{,}677$ & $0{,}662$ & $0{,}575$ & $0{,}495$ & $1/2$ \\
SVM (RBF) & $0{,}612$ & $0{,}599$ & $0{,}517$ & $0{,}469$ & $1/2$ \\
MLP & $0{,}539$ & $0{,}525$ & $0{,}453$ & $0{,}411$ & $1/2$ \\
\bottomrule
\end{tabular}

\end{table}

\section{Zuordnung und Reproduktion}
\label{sec:anh_reproduktion}

Tabelle~\ref{tab:anh_zuordnung} ordnet jede Abbildung, jede Tabelle und den Algorithmus dem erzeugenden oder prüfenden Skript und der Datengrundlage zu. \texttt{anhang\_\allowbreak tabellen.py} bricht ab, wenn ein Element der Kapitel~\ref{chap:einleitung} bis~\ref{chap:fazit} in der Zuordnung fehlt oder ein genanntes Skript nicht existiert. \texttt{nachweis\_\allowbreak zahlen.py} rechnet jede abgeleitete Zahl in Text und Tabellen aus den CSV-Dateien, der Konfiguration im Code und den gespeicherten Modellen nach und vergleicht sie mit dem LaTeX-Quelltext. Gerundet wird genau einmal. Nennt der Text die Differenz zweier gerundeter Werte, muss sie sich auch aus den gedruckten Werten ergeben.

\begin{table}[htbp]
	\centering
	\small
	\setlength{\tabcolsep}{3pt}
	\renewcommand{\arraystretch}{1.12}
	\caption[Zuordnung der Ergebnisse zu Skripten und Daten]{Zuordnung jeder Abbildung, jeder Tabelle und des Algorithmus zum erzeugenden oder prüfenden Skript und zur Datengrundlage.}
	\label{tab:anh_zuordnung}
	% erzeugt von anhang_tabellen.py, nicht von Hand ändern
\begin{tabular}{@{}L{2.8cm}L{4.6cm}L{7.4cm}@{}}
\toprule
\textbf{Element} & \textbf{Inhalt} & \textbf{Skript und Datengrundlage} \\
\midrule
Abbildungen~\ref{fig:paketstrom}, \ref{fig:generator_ablauf}, \ref{fig:blockschaltbild}, \ref{fig:qnet_architektur}, \ref{fig:backup} und~\ref{fig:eval_protokoll} & Schemata, als TikZ im Quelltext gesetzt & keine Messdaten \\
Tabellen~\ref{tab:rw_robustheit} und~\ref{tab:ki_vergleich} & Ergebnisse der Literatur & Quellen im Literaturverzeichnis \\
Abbildung~\ref{fig:feature_raum_3d} & Flüsse im Teilraum $(y_1,y_2,y_4)$ & \texttt{create\_3d\_plot.py}\newline Daten: Generator in \texttt{train\_eval.py} \\
Tabelle~\ref{tab:klassenueberlappung} & Abstand $D$ und Bayes-Fehler je SNR-Stufe & \texttt{nachweis\_zahlen.py} (Prüfung)\newline Daten: \texttt{bayes\_referenz.py} \\
Abbildung~\ref{fig:llr_signalraum} & Klassendichten der Prüfgröße $\Lambda$ & \texttt{llr\_signalraum\_plot.py}\newline Daten: \texttt{bayes\_referenz.py} \\
Algorithmus~\ref{alg:training}, Tabellen~\ref{tab:dqn_hyperparameters} und~\ref{tab:system_parameters} & Trainingsschleife und Konfiguration & \texttt{train\_eval.py}, \texttt{nachweis\_zahlen.py} (Prüfung)\newline Daten: Klasse \texttt{Config} \\
Tabelle~\ref{tab:ddqn_snr}, Abbildung~\ref{fig:snr_robustheit} & DDQN je SNR-Stufe & \texttt{train\_eval.py}, \texttt{generate\_plot.py}\newline Daten: \texttt{evaluation\_results.csv} \\
Abbildung~\ref{fig:entscheidungsfunktion} & Entscheidungsfunktion gegenüber der Bayes-Posterior & \texttt{entscheidungsfunktion\_plot.py}\newline Daten: \texttt{modelle/}, \texttt{baseline\_results.csv} \\
Tabelle~\ref{tab:f1_vergleich}, Abbildung~\ref{fig:baseline_vergleich} & F1-Score aller Verfahren & \texttt{eval\_baselines.py}, \texttt{generate\_plot.py}\newline Daten: beide CSV-Dateien \\
Tabelle~\ref{tab:reward} & mittlerer Reward je Fluss & \texttt{reward\_tabelle.py}\newline Daten: \texttt{modelle/}, beide CSV-Dateien \\
Tabelle~\ref{tab:rla_ela} & RLA und ELA & \texttt{nachweis\_zahlen.py} (Prüfung)\newline Daten: Tabelle~\ref{tab:f1_vergleich} \\
Tabelle~\ref{tab:bayes_grenze}, Abbildung~\ref{fig:pr_ebene} & Bayes-Referenz und gemessene Arbeitspunkte & \texttt{pr\_ebene\_plot.py}\newline Daten: \texttt{bayes\_referenz.py}, beide CSV-Dateien \\
Abbildung~\ref{fig:entscheidungsgrenzen} & Schnitt durch den Merkmalsraum & \texttt{entscheidungsgrenzen\_plot.py}\newline Daten: \texttt{modelle/ddqn\_seed0.pt}, \texttt{baseline\_results.csv} \\
Tabellen~\ref{tab:anh_prf1}, \ref{tab:anh_seeds}, \ref{tab:anh_schwellen} und~\ref{tab:anh_zuordnung} & Messwerte, Schwellen und Zuordnung & \texttt{anhang\_tabellen.py}\newline Daten: beide CSV-Dateien \\
Fließtext aller Kapitel & abgeleitete Zahlen in Text und Tabellen & \texttt{nachweis\_zahlen.py}\newline Daten: beide CSV-Dateien, Code, \texttt{modelle/} \\
Abschnitt~\ref{sec:reward_recall} & Kontrolllauf mit quadratischem Verlust & \texttt{kontrolle\_mse\_verlust.py}\newline Daten: eigener Lauf mit Seed~0 \\
Gewichte in \texttt{modelle/} & DDQN-Agenten der Seeds 0, 1 und 2 & \texttt{modelle\_nachtrainieren.py}\newline Daten: reproduziert \texttt{evaluation\_results.csv} \\
\bottomrule
\end{tabular}

\end{table}

Die Ergebnisse entstehen in dieser Reihenfolge:
\begin{enumerate}\raggedright
	\item \texttt{train\_\allowbreak eval.py}: Training und Evaluation des DDQN für die Seeds 0, 1 und 2, Ausgabe \texttt{evaluation\_\allowbreak results.csv}.
	\item \texttt{eval\_\allowbreak baselines.py}: Random Forest, SVM und MLP auf denselben Testflüssen, Ausgabe \texttt{baseline\_\allowbreak results.csv}.
	\item \texttt{modelle\_\allowbreak nachtrainieren.py}: trainiert die drei DDQN-Agenten mit identischen Seeds erneut und speichert die Gewichte in \texttt{modelle/}.
	\item Plotskripte aus Tabelle~\ref{tab:anh_zuordnung} für die Abbildungen.
	\item \texttt{reward\_\allowbreak tabelle.py} und \texttt{anhang\_\allowbreak tabellen.py} für Tabelle~\ref{tab:reward} und die Tabellen dieses Anhangs.
	\item \texttt{nachweis\_\allowbreak zahlen.py} als abschließende Prüfung, Exit-Code~1 bei jeder Abweichung.
\end{enumerate}

Gerechnet wurde mit Python~3.14.5, PyTorch~2.11.0 (CUDA~12.8), scikit-learn~1.9.1, NumPy~2.5.4, SciPy~1.18.1, Gymnasium~1.4.0 und Matplotlib~3.11.2 auf einer NVIDIA GeForce RTX~2060. In dieser Umgebung reproduziert das Nachtraining alle 36 DDQN-Werte der \texttt{evaluation\_\allowbreak results.csv} auf vier Stellen, und die neu trainierten Baselines liefern dieselben Werte wie \texttt{baseline\_\allowbreak results.csv}. Skripte, CSV-Dateien und Modellgewichte gehören zum Projektarchiv der Arbeit.
